%%%%%%%%%%%%%%%%%%%%%%%%%%%%%%%%%%%%%%%%%
% Stefano Cover Letter
% LaTeX Template
% Version 2.0 (February 14, 2023)
%
% This template originates from:
% https://www.LaTeXTemplates.com
%
% Authors:
% Vel (vel@latextemplates.com)
% Stefano (https://www.kindoblue.nl)
%
% License:
% CC BY-NC-SA 4.0 (https://creativecommons.org/licenses/by-nc-sa/4.0/)
%
%%%%%%%%%%%%%%%%%%%%%%%%%%%%%%%%%%%%%%%%%

%----------------------------------------------------------------------------------------
%	PACKAGES AND OTHER DOCUMENT CONFIGURATIONS
%----------------------------------------------------------------------------------------

\documentclass[
	%letterpaper, % Uncomment for US letter paper size
	parskip=half, % Vertical whitespace between paragraphs
	%foldmarks=hp, % Uncomment to remove fold marks on the left side of the first page
	enlargefirstpage=true, % Allow more text on the first page
]{scrlttr2} % Use the KOMA-Script letter class

\usepackage{fontsize} % Allows changing the default font size to an arbitrary value
\changefontsize[16pt]{12pt} % Change the default font size and baselineskip
\usepackage{anyfontsize} % Allows font sizes at arbitrary values

\setlength{\parindent}{0pt} % Paragraph indentation

%\pagestyle{empty} % Uncomment this for no footer on the second and onwards pages

%----------------------------------------------------------------------------------------
%	FONTS
%----------------------------------------------------------------------------------------

\usepackage[utf8]{inputenc} % Required for inputting international characters
\usepackage[T1]{fontenc} % Output font encoding for international characters

\usepackage{ebgaramond} % Use the EB Garamond font as the main serif font
\usepackage{lato} % Use the Lato font as the sans font

\usepackage[letterspace=50]{microtype} % Required for changing letter spacing

\usepackage{fontawesome5} % Allows the use of Font Awesome icons
\usepackage[hidelinks]{hyperref}

%----------------------------------------------------------------------------------------
%	COVER LETTER INFORMATION
%----------------------------------------------------------------------------------------

\setkomavar{fromname}{Yiran Liu} % Your name
\setkomavar{signature}{Yiran Liu} % Your name as you want it to appear in the signature

\newkomavar{subtitle}\setkomavar{subtitle}{Computer Engineering Student} % Set the subtitle line for the first page header, leave empty for no subtitle

\setkomavar{fromaddress}{} % Your address or location
\setkomavar{fromphone}{+1 (437) 535-8276} % Your phone number
\setkomavar{fromemail}{\href{mailto:y3275liu@uwaterloo.ca}{y3275liu@uwaterloo.ca}} % Your email address

\setkomavar{place}{} % A city or country that's output before the date automatically, add your location here if you want this
\setkomavar{date}{\today} % Date output above the letter, use \today for today's date or leave this empty for no date

\setkomavar{subject}{Application for Robotics Software Developer Co-op Position} % Bold text above the letter specifying what it's for, useful for specifying the job you are applying for. Leave this empty for no subject.

% These are not used in this document, uncomment them if you would like to use them and output their values with e.g.: \usekomavar{fromurl}
%\setkomavar{fromfax}{+1 (123) 456 7890} % Your fax number
\setkomavar{fromurl}{\href{https://github.com/catr1xLiu/}{github.com/catr1xLiu/}} % Your personal website

%------------------------------------------------

\setkomavar{firsthead}{ % Style the first page header
	\centering\scshape % Text position and styling
	{\huge\textls{\usekomavar{fromname}}}\\[4pt] % Your name
	{\large\usekomavar{subtitle}} % Subtitle
}

\setplength{firstheadvpos}{0.045\paperheight} % Vertical position of the header on the first page

%------------------------------------------------

\setkomavar{firstfoot}{ % Style the first page footer
	\centering\normalfont % Text position and styling
	{\small{\footnotesize\faEnvelope}\hspace{4pt}\usekomavar{fromemail}\hspace{8pt}{\footnotesize\faMobile*}\hspace{4pt}\usekomavar{fromphone}\hspace{8pt}{\footnotesize\faGithub}\hspace{4pt}\usekomavar{fromurl}}
}

\setplength{firstfootvpos}{0.96\paperheight} % Vertical position of the footer on the first page
\setplength{refvpos}{7.5cm} % Move letter body up to reduce white space (DIN default is ~9cm)

%----------------------------------------------------------------------------------------

\begin{document}

%----------------------------------------------------------------------------------------
%	COVER LETTER CONTENT
%----------------------------------------------------------------------------------------

\begin{letter}{
	Red Rabbit AI Canada ULC \\
	Divisional Office \\
	6378 Silver Avenue \\
	Burnaby, BC V5H 0J2 \\
	Canada
}

\opening{Dear Hiring Manager,}

I'm a 1B Computer Engineering student at UWaterloo applying for the 4-month Fall 2026 Robotics Software Developer Co-op. Red Rabbit's hybrid teleop and AI approach to humanoids aligns perfectly with my recent research on creating real-time teleoperation pipelines that support learned-policy training.

At the National Research Council (NRC) Canada, I designed a force-aware teleoperation system that allows operators to feel what the robot feels during data collection. This used a 6-DoF Haply Inverse3 haptic controller, a UR10e arm, and wrist-mounted force-and-torque sensors. I improved the policy's inference pipeline by making it modular with a three-layer architecture and by adding Real-Time Chunking to keep the robot's motion smooth, even with machine-learning model delays. In our RoboTac 2026 submission, the end-to-end system reduced peak contact force in demonstrations by 64\% compared to the force-blind baseline.

I'm a practical builder, not just a researcher. Before NRC, I created MapleSim, an open-source rigid-body physics library used by over 500 FIRST Robotics teams worldwide and received more than 100 GitHub stars. During FRC, I also developed a robot localization system that accurately determines position on the field within centimeters, utilizing AprilTag fiducial markers combined with wheel odometry, which I debugged with log-replay simulation. I spent several competition seasons fixing mechanical and electrical issues on the robot. I'm skilled in Python and C++, and my experience with WPILib (Java) control code directly relates to ROS2 patterns.

For your specific environment: I'm a native Mandarin speaker. I'm also proficient in agentic engineering, regularly customizing Claude Code and OpenCode (skills, restrictions, MCP servers, inference backends) to maximize efficiency.

I would love to bring my teleoperation experience and hardware troubleshooting skills to your humanoid mission this Fall. Thank you for your time, and I look forward to your response.

Best, \\
\usekomavar{signature}

%----------------------------------------------------------------------------------------

\end{letter}

\end{document}
